\scanpage{33}

% L1–2
$V$ groupe comm. (écrit additivement)

$I$ ens de cardinal 3

% L3
\[
 \Gamma=\mathrm{Ker}\,(V^I\to V)
\]

% L4
$\omega$ ens des \inserted{\unsure{2|?}} ordres circulaires sur $I$ (c'est un $(\pm1)$-\unsure{tors|tors.})

% L5
\[
 V^\omega=V\mathbin{\wedge}_{\{\pm1\}}\omega
 \qquad(\pm1\text{ opérant sur }V\text{ par chg}\textsuperscript{t}\text{ de signe})
\]

% L6–9
Pour $i\in I$, \unsure{---}
\[
 V^\omega\xrightarrow{\varphi_i}\Gamma\quad\text{par}\quad\deleted{(x,\rho)}
 \left\{\begin{aligned}
 &\varphi_i(x\wedge\rho)_{\deleted{\unsure{\rho}}i}=0\\
 &\varphi_i(x\wedge\rho)_{\rho i}=x\\
 &\varphi_i(x\wedge\rho)_{\rho^2 i}=-x
 \end{aligned}\right.
\]

% L6–10 right
\margin{\uline{NB} \(\deleted{\unsure{\varphi_i(x)+\varphi_{\rho i}(x)+\varphi_{\rho^2 i}(x)=0}}\)
\(\displaystyle\sum_{i\in I}\varphi_i(x)=0\quad\forall x\in V^\omega\)}

% L10
On cherche à obtenir les systèmes

% L11–13
\deleted{\unsure{et|$\psi$}}
\diagram{\begin{tikzcd}[ampersand replacement=\&]
 V^\omega\arrow[r,dash,"\varphi_i"]\arrow[dr,"\psi_i"']\&\Gamma\\
 \&\widetilde\Gamma^{\unsure{''|*}}\arrow[u,"{\unsure{u|k}}"']
\end{tikzcd}}

% L10–14 right
\margin{\deleted{\uline{NB}} En fait, on a suite exacte
\(\begin{gathered}0-V^\omega\xrightarrow[\mathrm{diag}]{}(V^\omega)^I\to\Gamma-0\\
\text{i.e.}\\
0-V^\omega-(V^\omega)^I\to V^I\to V\to0\end{gathered}\)}

% L14–17
avec

\deleted{\unsure{i}} a) $\widetilde\Gamma\xrightarrow{u}\Gamma$ est un hom de groupes à noyau central $Z$

b) $\psi_i$ un hom $V^\omega\mathbin{\unsure{-|\to}}\widetilde\Gamma$ qui relève $\varphi_i$ :

% L18
\[
 u\psi_i=\varphi_i
\]

% L19–20
\uline{NB} $u$ doit être épi, donc $\widetilde\Gamma$ est une \uline{extension} de $\Gamma$ par $Z$ (écrit addit\textsuperscript{\unsure{mt}})

% L21
Soit $\rho\in\omega$, $i_0\in I$, d'où $i_1=\rho i_0,\ i_2=\rho^2i_0=\rho i_1$. On aura

% L22
\[
 \varphi_{i_0}(\deleted{\unsure{X}}x\wedge\rho)+\varphi_{i_1}(x\wedge\rho)+\varphi_{i_2}(x\wedge\rho)=0
\]

% L23
d'où
\[
 \psi_{i_0}(x\wedge\rho)\,\psi_{i_1}(x\wedge\rho)\,\psi_{i_2}(x\wedge\rho)
 =z(x,\rho,i_0)\in Z.
\]
% archive sheet number: 32
